% Generated draft: requires mathematical and source review.
\documentclass[11pt]{article}
\usepackage[margin=25mm]{geometry}
\usepackage{fontspec}
\setmainfont{Latin Modern Roman}
\usepackage{amsmath,amssymb,mathtools,unicode-math}
\setmathfont{Latin Modern Math}
\usepackage{enumitem,needspace,xurl,hyperref}
\setlength{\parindent}{0pt}
\setlength{\parskip}{5pt}
\setlength{\emergencystretch}{3em}
\newcommand{\N}{\mathbb{N}}
\newcommand{\Z}{\mathbb{Z}}
\newcommand{\Q}{\mathbb{Q}}
\newcommand{\R}{\mathbb{R}}
\newcommand{\C}{\mathbb{C}}
\newcommand{\F}{\mathbb{F}}
\newcommand{\A}{\mathbb{A}}
\newcommand{\PP}{\mathbb{P}}
\newcommand{\E}{\mathbb{E}}
\newcommand{\eps}{\varepsilon}
\newcommand{\dd}{\,\mathrm{d}}
\begin{document}
\section*{UnsolvedMath ID: 11100035}
\textbf{Source problem code:} catalogue entry
\subsection{Definitions and mathematical statement}
Fix the stable homotopy category $\mathcal T$ in the source's axiomatic setting. The exact axioms, size universe, and intended scope beyond spectra are not reproduced. A localizing subcategory is a full replete triangulated subcategory closed under arbitrary coproducts; a colocalizing subcategory is similarly closed under arbitrary products. Replete means closed under isomorphism in the ambient category.

\textbf{Axiomatic stable homotopy problem 3 (Hovey).} Prove that localizing subcategories of $\mathcal T$ form only a set, rather than a proper class. Concretely, seek a set $I$ and localizing subcategories $\mathcal L_i$ for $i\in I$ such that every full replete localizing subcategory of $\mathcal T$ equals some $\mathcal L_i$. The target concerns subcategories of the fixed ambient category, not merely abstract equivalence types, and the source does not restrict to a preselected set-generated family.

In a monoidal setting, a Bousfield class of $E$ is determined by its acyclic objects,
\[
 \langle E\rangle=\{X\in\mathcal T:E\otimes X\cong0\}.
\]
The source separately reports that there are only set-many Bousfield classes, attributing this to Ohkawa with simplified proofs by Strickland and Dwyer--Palmieri. It expects a one-to-one correspondence among Bousfield classes, localizing subcategories, and colocalizing subcategories, while explicitly saying that this correspondence may be out of reach.

The established Bousfield set-sizedness report is not itself the requested localizing-subcategory theorem. An injection or classification of arbitrary localizing subcategories by the set of Bousfield classes would need proof; a passage to orthogonals does not automatically supply such a classification under unspecified hypotheses. The registry summary claims the broader result is solved by this consequence, but provides no precise theorem establishing that missing step and partly restates the source's expectation as though already known. That argument is not adopted as a solution.
\subsection{Short English statement}
Prove that a fixed stable homotopy category has only set-many localizing subcategories. The known set-sizedness of Bousfield classes and the hoped-for correspondence with localizing and colocalizing subcategories are distinct statements.
\subsection{Sources}
\begin{itemize}
\item Local UnsolvedMath registry, permanent ID 11100035, source code catalogue entry; attributed to Mark Hovey, no proposal year supplied.
\item \emph{Hovey's Algebraic Topology Problem List}, Axiomatic stable homotopy, item 3; surviving mirror \url{https://www-users.cse.umn.edu/~tlawson/hovey/}. The registry records HTML-ordered-list extraction and access on 2026-07-29; the mirror has 65 top-level items on nine category pages.
\item Ohkawa, Strickland, and Dwyer--Palmieri are named for Bousfield-class set-sizedness/proofs. The summary also mentions a Hovey--Palmieri--Strickland memoir, but gives no complete bibliography or precise localizing-subcategory theorem. Its remaining bibliographic fragments and orthogonal argument are incomplete and unverified.
\end{itemize}
\end{document}
